\begin{tikzpicture}[scale=0.92, transform shape, >=stealth]

    % ========================================================
    % 1. GENERADOR DE FUNCIONES (GW-Instek GFG-3015)
    % ========================================================
    \draw[thick, fill=gray!8] (-0.6, 0.2) rectangle (3.2, 5.0);
    \node[font=\bfseries\small, align=center] at (1.2, 2.7) {Generador de\\Funciones\\[5pt]\footnotesize GW-Instek\\GFG-3015};
    
    % Conectores BNC panel frontal (x = 3.2)
    \coordinate (gen_out)  at (3.2, 4.2);
    \coordinate (gen_sync) at (3.2, 1.2);
    \coordinate (gen_gnd)  at (3.2, 0.2);
    
    \draw[fill=white, thick] (gen_out) circle (3.5pt);
    \node[left=4pt, font=\scriptsize\bfseries] at (gen_out) {MAIN OUT};
    
    \draw[fill=white, thick] (gen_sync) circle (3.5pt);
    \node[left=4pt, font=\scriptsize\bfseries] at (gen_sync) {SYNC OUT};

    \draw (1.3, 0.2) node[ground]{};

    % ========================================================
    % 2. OSCILOSCOPIO DIGITAL (Tektronix TDS 1001)
    % ========================================================
    \draw[thick, fill=gray!8] (11.0, 0.2) rectangle (14.8, 5.0);
    \node[font=\bfseries\small, align=center] at (13.0, 2.7) {Osciloscopio\\Digital\\[5pt]\footnotesize Tektronix\\TDS 1001};

    % Conectores BNC panel frontal (x = 11.0)
    \coordinate (osc_ch1)  at (11.0, 4.2);
    \coordinate (osc_ch2)  at (11.0, 2.1);
    \coordinate (osc_trig) at (11.0, 1.2);
    \coordinate (osc_gnd)  at (11.0, 0.2);
    
    \draw[fill=white, thick] (osc_ch1) circle (3.5pt);
    \node[right=4pt, font=\scriptsize\bfseries, color=blue!80!black] at (osc_ch1) {CH1};
    
    \draw[fill=white, thick] (osc_ch2) circle (3.5pt);
    \node[right=4pt, font=\scriptsize\bfseries, color=red!80!black] at (osc_ch2) {CH2};
    
    \draw[fill=white, thick] (osc_trig) circle (3.5pt);
    \node[right=4pt, font=\scriptsize\bfseries] at (osc_trig) {EXT TRIG};

    \draw (12.9, 0.2) node[ground]{};

    % ========================================================
    % 3. BANCO DE PRUEBA RLC (Placa MEI)
    % ========================================================
    \def\xBranch{6.8}
    
    % Resistor serie de protección r = 47 ohm (con separación holgada hacia Nodo A)
    \draw (gen_out) to[short] (4.0, 4.2)
          to[R, l={$r = \qty{47}{\ohm}$}, *-*] (5.8, 4.2)
          -- (\xBranch, 4.2) coordinate (nodeA);
    \node[above=4pt, font=\scriptsize\bfseries, color=blue!85!black] at (nodeA) {Nodo A};

    % Rama resonante vertical: Inductor L + Resistencia interna RL
    \draw (nodeA) to[L, l_={$L = \qty{30}{\milli\henry}$}] (\xBranch, 3.2)
          to[R, l_={$R_L$}] (\xBranch, 2.3)
          -- (\xBranch, 2.1) coordinate (nodeB);
    \node[left=5pt, font=\scriptsize\bfseries, color=red!85!black] at (nodeB) {Nodo B};

    % Capacitor C a masa
    \draw (nodeB) to[C, l_={$C = \qty{7}{\nano\farad}$}, *-*] (\xBranch, 0.9) coordinate (nodeGndRLC);
    
    % Conexión a la línea de masa común
    \draw (nodeGndRLC) -- (\xBranch, 0.2);
    \draw (\xBranch, 0.2) node[ground]{};

    % ========================================================
    % 4. CONEXIÓN DE PUNTAS Y REFERENCIAS (CH1, CH2, DISPARO)
    % ========================================================
    
    % CANAL 1 (CH1): Mide e_tot en Nodo A
    \draw[blue!85!black, very thick] (nodeA) -- (osc_ch1);
    \node[above=3pt, font=\footnotesize\bfseries, color=blue!85!black] 
        at ($(\xBranch, 4.2)!0.55!(osc_ch1)$) {$e_{\text{tot}}$ (Canal 1)};
    \draw[blue!85!black, fill=blue!85!black] (nodeA) circle (2.5pt);

    % CANAL 2 (CH2): Mide e_c en Nodo B
    \draw[red!85!black, very thick] (nodeB) -- (osc_ch2);
    \node[above=3pt, font=\footnotesize\bfseries, color=red!85!black] 
        at ($(\xBranch, 2.1)!0.55!(osc_ch2)$) {$e_c$ (Canal 2)};
    \draw[red!85!black, fill=red!85!black] (nodeB) circle (2.5pt);

    % LÍNEA DE MASA COMÚN
    \draw[black, thick] (gen_gnd) -- (osc_gnd);

    % DISPARO EXTERNO (SYNC -> EXT TRIG): Por debajo de la masa
    \draw[gray!50!black, thick, dashed] (gen_sync) -- ++(0.4, 0) -- ++(0, -1.8)
          -- ($(osc_trig) + (-0.4, -1.8)$)
          node[midway, below=2pt, font=\scriptsize\bfseries] {Sincronismo TTL (\emph{Sync Out} $\to$ \emph{Ext Trig})}
          -- ++(0, 1.8) -- (osc_trig);

\end{tikzpicture}
